\title{Direct Preference Optimization: Your Language Model is Secretly a Reward Model}

\begin{abstract}
Large-scale unsupervised language models (LMs) acquire extensive general knowledge and some degree of reasoning capabilities; however, exerting precise influence over their outputs remains challenging because their training lacks supervision. To address this, current approaches gather human judgments comparing the quality of generated text and adapt the LMs to prefer these favored responses, typically using reinforcement learning from human feedback (RLHF). Yet, RLHF is inherently intricate and often prone to instability: it first involves training a reward model to encode human preferences and subsequently fine-tuning the original LM through reinforcement learning to optimize this reward signal, all while constraining the model to remain close to its initial state. In this work, we propose an alternative reward model formulation within the RLHF paradigm that allows for direct, closed-form derivation of the associated optimal policy. This reformulation transforms the traditional RLHF objective into a simple classification problem, which we term \textit{Direct Preference Optimization} (DPO). DPO offers a stable, efficient, and computationally simple solution, dispensing with the need to sample from the language model during fine-tuning and minimizing reliance on hyperparameter adjustments. Empirically, we demonstrate that DPO can adapt LMs to mirror human preferences at least as effectively as, if not better than, prevalent strategies. Remarkably, DPO surpasses PPO-based RLHF in tasks demanding sentiment control, and achieves comparable or superior performance in summarization and single-turn dialogue tasks, all while greatly simplifying the fine-tuning process.
\end{abstract}

\section{Introduction}
Large-scale unsupervised language models (LMs), trained on immense corpora, have demonstrated unexpected and impressive capabilities~\citep{chowdhery2022palm, brown2020language, touvron2023llama,bubeck2023sparks}. Yet, the training data for these models is generated by humans, each with differing intentions, levels of expertise, and objectives. As a result, not every human skill or intent in the dataset should be reproduced; for instance, it is beneficial for an AI programming assistant to \textit{recognize} frequent programming errors for correction purposes, but, in code generation tasks, the goal is to favor the (potentially uncommon) high-caliber coding present in some portions of the dataset. In a similar vein, a language model should ideally be \textit{informed} of prevalent misconceptions—such as a false belief held by 50\% of people—without replicating the error in its responses half the time when queried about it. Put differently, to design AI systems that are reliable, high-performing, and manageable, one must carefully control the model's \emph{preferred outputs and conduct} amidst its extensive repertoire of \textit{knowledge and skills} \citep{ouyang2022training}. Presently, mainstream approaches primarily rely on reinforcement learning (RL) techniques to align LMs with human preferences, but we argue that the RL-inspired objectives often used in such methods are exactly achievable via a straightforward binary cross-entropy loss, thereby streamlining the overall preference alignment process.

In general, prevailing techniques shape desired LM behavior through the use of carefully selected datasets that reflect human judgments of helpful and safe actions. This process, referred to as preference learning, is usually performed subsequent to comprehensive unsupervised pre-training on large text corpora. Although the simplest way to incorporate preferences is to use supervised fine-tuning on top-tier human demonstrations, methods based on reinforcement learning from human or artificial feedback (RLHF/RLAIF; \citep{christiano2017deep,bai2022constitutional}) have proven most effective. RLHF employs a learned reward model, which is trained on human preference datasets, to guide an RL-based policy that generates responses with higher reward scores, all the while minimizing divergence from the pre-trained model. Despite RLHF’s successes in conversational and programming domains, its training regimen is substantially more complicated than standard supervised learning: it involves multiple language models, policy sampling during training, and a much greater computational burden.

This work introduces a method to optimize language models directly for human preferences, foregoing explicit reward modeling and reinforcement learning techniques. We present \textit{Direct Preference Optimization (DPO)}, an approach that implicitly targets the same objective as established RLHF algorithms—maximizing reward subject to a KL-divergence penalty—yet remains easy to implement and train. Conceptually, DPO updates function by amplifying the log probability of preferred outputs relative to less-preferred ones. Crucially, these updates involve per-sample importance weights, which help mitigate the degeneration observed when using a simplistic probability ratio objective. Similar to existing approaches, DPO adopts a theoretical framework for preference modeling (e.g., the Bradley-Terry model; \cite{bradley1952rankanalysis}), which quantifies the agreement between a reward function and observed preference data. Unlike conventional methods that apply this preference model to train a reward function and subsequently optimize a policy with respect to this learned reward, DPO reparameterizes the problem, formulating the preference loss as a direct function of the policy. Using human-labeled preference data over generated responses, DPO can thus optimize policies via a binary cross-entropy objective, arriving at an optimal policy that implicitly aligns with the reward signal inferred from the preferences.

The principal contribution of this study is Direct Preference Optimization (DPO), an algorithm for preference-based language model training that eliminates the need for reinforcement learning. Empirical evaluations reveal that DPO performs on par with, or better than, leading alternatives such as PPO-based RLHF across tasks including sentiment control, summarization, and dialogue, utilizing language models containing up to 6B parameters.

\section{Related Work}

Self-supervised language models, when scaled up, have demonstrated the capability to perform certain tasks in a zero-shot manner \citep{radford2019language} or by utilizing few-shot prompting \citep{gpt3,megatron,chowdhery2022palm}. Nevertheless, their effectiveness on downstream applications and their ability to conform to user intent are often substantially enhanced through fine-tuning on datasets comprised of instructions and human-provided completions \citep{mishra-etal-2022-cross,sanh2022multitask,chung2022scaling,thoppilan2022lamda}. This process, termed ‘instruction-tuning’, allows large language models (LLMs) to extrapolate to new instructions beyond the original set and generally leads to improvements in practical usability \citep{chung2022scaling}. Although instruction tuning has been notably successful, collecting \textit{comparative} human evaluations of response quality is frequently more scalable than assembling expert-generated demonstrations. As a result, follow-up research has focused on fine-tuning LLMs using human preference datasets, which has enhanced their capabilities in translation \citep{kreutzer-etal-2018-reliability}, summarization \citep{stiennon2022learning,ziegler2020finetuning}, creative writing \citep{ziegler2020finetuning}, and following instructions \citep{ouyang2022training,ramamurthy2023is}. In these approaches, a neural reward model is first trained to fit the human preferences dataset using frameworks such as the Bradley-Terry model \citep{bradley1952rankanalysis}. Subsequently, the LLM is fine-tuned via reinforcement learning algorithms—including REINFORCE \citep{williams1992reinforce}, proximal policy optimization (PPO; \cite{schulman2017proximal}), or their derivatives \citep{ramamurthy2023is}—to maximize the learned reward. Parallel research has taken advantage of LLMs refined for instruction adherence with human feedback, harnessing them to create synthetic preference datasets targeting characteristics like safety or non-harmfulness \citep{bai2022constitutional}; this is achieved by relying on weak human supervision through text-based rubrics for LLM assessments. Collectively, these advancements bridge two research domains: one centered on applying reinforcement learning to language model training for diverse objectives~\citep{Ranzato2015SequenceLT,paulus2018a,wu2018learning}, and another focused on learning broadly from human preferences \citep{christiano2017deep,kupcsik2018learning}. Although leveraging relative human preferences is compelling, reinforcement learning-based fine-tuning of large language models presents substantial practical difficulties; the present work introduces a theoretically-grounded strategy for optimizing relative preferences that obviates the need for reinforcement learning.

Beyond applications involving language, policy learning from preferences has been explored within both the bandit and reinforcement learning paradigms, leading to a variety of methodologies. The contextual dueling bandit (CDB; \cite{yue2012karmed, dudik2015contextual}) framework specifically addresses situations where learning is driven by preferences or rankings over actions instead of explicit rewards. In contexts where absolute reward signals are unavailable, theoretical treatments of CDBs replace the classic optimal policy with the concept of a \textit{von Neumann winner}—that is, a policy whose probability of outperforming any alternative policy meets or exceeds 50\% \citep{dudik2015contextual}. Notably, CDB setups typically assume preference labels are provided in an online manner. By contrast, in the domain of learning from human preferences, learning often proceeds from a static, offline collection of preference-annotated action pairs \citep{yan2022human}. Similarly, the area of \textit{preference-based RL} (PbRL) focuses on learning using binary feedback generated by an \textit{unknown} scoring function, rather than direct rewards \citep{BusaFekete2014, ruiz2023dueling}. A range of PbRL algorithms exist, some capable of utilizing off-policy preference information, but these approaches usually involve first estimating the hidden scoring (or reward) function before optimizing against it \citep{jain2013learning, BusaFekete2014, christiano2017deep, sadigh2017active, kupcsik2018learning}. By contrast, the method presented here advocates for a unified approach that learns a policy in a single stage by directly aligning with observed preferences.

\section{Preliminaries}\label{section:prelims}

We now provide a summary of the RLHF framework as delineated by \citeauthor{ziegler2020finetuning} (see also \citep{stiennon2022learning, bai2022training, ouyang2022training}), which typically comprises three distinct phases: (1) supervised fine-tuning (SFT); (2) preference data collection coupled with reward model training; and (3) reinforcement learning-based optimization.

\textbf{SFT}: The RLHF workflow customarily begins by adapting a pre-trained language model to the target domain via supervised learning on carefully curated data relevant to the downstream application (such as dialogue or summarization), resulting in the fine-tuned model $\pi^\text{SFT}$.

\textbf{Reward Modelling Phase}: During the subsequent stage, the SFT model is prompted with various inputs $x$ to generate answer pairs $(y_1, y_2)\sim \pi^\text{SFT}(y \mid x)$. These response pairs are then assessed by human evaluators, who indicate a preferred answer—denoted as $y_w\succ y_l \mid x$, with $y_w$ and $y_l$ representing the chosen and rejected completions, respectively. Although the underlying generator of these preferences is a latent reward function $r^*(y, x)$, this function is not directly observable. Preferences may be modelled in multiple ways; a commonly utilized approach is the Bradley-Terry (BT) model \cite{bradley1952rankanalysis}, though generalizations such as the Plackett-Luce model \citep{plackett1975analysis, luce2012individual} are also compatible if ranked outputs are available. The BT framework prescribes that the distribution over human preferences $p^*$ can be formulated as:

Given a fixed dataset of comparison examples $\mathcal{D}=\bigl\{x^{(i)}, y_w^{(i)}, y_l^{(i)}\bigr\}_{i=1}^N$ sampled according to $p^*$, we introduce a reward model parameterized as $r_{\phi}(x, y)$, whose parameters can be learned via maximum likelihood estimation. When framed as a binary classification problem, this leads to the negative log-likelihood loss function:

\begin{equation}\label{eq:reward_model}
    \mathcal{L}_R(r_{\phi}, \mathcal{D}) = -\mathbb{E}_{(x, y_w, y_l)\sim \mathcal{D}}\bigl[\log \sigma(r_{\phi}(x, y_w)- r_{\phi}(x, y_l))\bigr]
\end{equation}

where $\sigma$ denotes the logistic sigmoid function. For large language models, $r_{\phi}(x, y)$ is typically initialized with the supervised fine-tuned (SFT) model $\pi^\text{SFT}(y \mid x)$, followed by appending a linear head to the output of the final transformer block, yielding a scalar reward estimate \cite{ziegler2020finetuning}. In order to decrease the variance of the learned reward, normalization is applied in earlier works so that  $\mathbb{E}_{x,y\sim \mathcal{D}}\left[r_\phi(x, y)\right] = 0$ for every $x$.

\textbf{RL Fine-Tuning Phase}: In the reinforcement learning phase, this estimated reward function acts as a feedback signal for guiding the language model's policy. Following established techniques~\citep{jaques2017sequence, jaques2020human}, the optimization problem takes the form

\begin{equation}\label{eq:RL}
\max_{\pi_{\theta}}  \mathbb{E}_{x\sim \mathcal{D}, y\sim \pi_{\theta}(y \mid x)}\bigl[r_{\phi}(x, y)\bigr] - \beta\mathbb{D}_{\textrm{KL}}\bigl[\pi_{\theta}(y\mid x)\mid \mid \pi_\text{ref}(y\mid x)\bigr],
\end{equation}

with the parameter $\beta$ modulating how much $\pi_{\theta}$ can deviate from the base policy $\pi_\text{ref}$, usually set as the initial SFT model $\pi^\text{SFT}$. In most implementations, the initial weights of $\pi_\theta$ are also copied from $\pi^\text{SFT}$. This regularization term is essential: it keeps the updated model within regions of policy space where the reward estimator remains reliable and prevents issues such as reduced diversity or collapse to a handful of high-scoring outputs. Given the nondifferentiable nature of language generation, this objective is commonly optimized via reinforcement learning. The prevailing method \citep{ziegler2020finetuning, stiennon2022learning, bai2022training, ouyang2022training} designs the reward for RL as ${r(x, y) = r_{\phi}(x, y) -\beta (\log \pi_{\theta}(y\mid x) - \log \pi_\text{ref}(y\mid x))}$, optimized with algorithms such as PPO \cite{schulman2017proximal}.

\section{Direct Preference Optimization}\label{sec:DPO}

With the complexity and inefficiencies inherent to scaling reinforcement learning approaches for language model fine-tuning, we aim to formulate a more straightforward policy optimization strategy based on preference data. In contrast to conventional RLHF pipelines that first construct a reward and subsequently apply RL to maximize it, our method makes use of a specific parameterization of the reward model that enables us to recover its optimal policy analytically—thus sidestepping the RL loop altogether. As detailed below, our principal innovation involves utilizing a closed-form relationship mapping reward functions to their optimal policies, which facilitates rewriting losses defined over rewards directly as losses over policies. By effecting this variable substitution, we dispense with explicit reward model training, while remaining consistent with established models for human preference, such as the Bradley-Terry framework. Thus, the policy network can simultaneously represent the generative language model and encode implicit rewards.

\textbf{DPO Objective Derivation.} We begin by considering the same reinforcement learning objective outlined in previous literature, specifically Eq.~\ref{eq:RL}, but allowing for a generic reward function $r$. Prior studies~\citep{peters2007reinforcement, peng2019advantage, korbak2022reinforcement, go2023aligning} have established that the optimal policy for the KL-regularized reward maximization problem (Eq.~\ref{eq:RL}) can be expressed as:

\begin{equation}\label{eq:op_policy}
    \pi_r(y\mid x) = \frac{1}{Z(x)}\pi_\text{ref}(y\mid x)\exp\left(\frac{1}{\beta}r(x, y)\right),
\end{equation}

where the normalization constant is defined by $Z(x) =\sum_{y}\pi_\text{ref}(y\mid x)\exp\left(\frac{1}{\beta}r(x, y)\right)$. A detailed derivation is available in Appendix \ref{app:derivation1}. In practice, direct estimation of the partition function $Z(x)$, even when working with the MLE-based reward approximation $r_\phi$ for $r^*$, remains computationally intensive~\citep{korbak2022reinforcement, go2023aligning}. To address this issue, Eq.~\ref{eq:op_policy} can be manipulated to represent the reward in terms of the induced optimal policy $\pi_r$, the baseline policy $\pi_\text{ref}$, and the (unobserved) normalization $Z(\cdot)$. Taking the natural logarithm of both sides of Eq.~\ref{eq:op_policy} and rearranging yields:

\begin{equation}\label{eq:main_eq}
    r(x,y) =\beta \log \frac{\pi_r(y\mid x)}{\pi_\text{ref}(y\mid x)} + \beta \log Z(x).
\end{equation}

This alternative parameterization is applicable for the true reward $r^*$ as well as its optimal policy $\pi^*$. The Bradley-Terry framework utilizes only the reward differences between two choices; specifically, $p^*(y_1 \succ y_2 \mid x) = \sigma(r^*(x, y_1) - r^*(x, y_2))$. If we insert the expression from Eq.~\ref{eq:main_eq} for $r^*(x,y)$ into the pairwise preference model Eq.~\ref{eq:bradley-terry}, the partition functions for $y_1$ and $y_2$ cancel, resulting in a human preference probability reliant solely on the optimal policy $\pi^*$ and the reference policy $\pi_\text{ref}$. Hence, for RLHF under the Bradley-Terry assumption, the optimal policy $\pi^*$ conforms to the following preference model:

\begin{equation}\label{eq:objective}
    p^*(y_1\succ y_2 \mid x)=\frac{1}{1 + \exp\left(\beta \log \frac{\pi^*(y_2\mid x)}{\pi_\text{ref}(y_2\mid x)} - \beta \log \frac{\pi^*(y_1\mid x)}{\pi_\text{ref}(y_1\mid x)}\right)}
\end{equation}

Appendix~\ref{app:derivation2} provides the detailed steps for this result. Although this is shown for the Bradley-Terry model, analogous expressions follow for the more general Plackett-Luce family~\citep{plackett1975analysis, luce2012individual}, as shown in Appendix~\ref{app:plackett_luce_models}.

With this reformulation—where the probability of human preferences is written as a function of the optimal policy rather than directly in terms of a reward function—we can derive a maximum likelihood training objective for a parametric policy $\pi_\theta$. In analogy to conventional reward modeling (see Eq.~\ref{eq:reward_model}), this leads to a policy objective given by:

\begin{equation}\label{eq:optimum_model}
    \mathcal{L}_\text{DPO}(\pi_{\theta}; \pi_\text{ref}) = -\mathbb{E}_{(x, y_w, y_l)\sim \mathcal{D}}\left[\log \sigma \left(\beta \log \frac{\pi_{\theta}(y_w\mid x)}{\pi_\text{ref}(y_w\mid x)} - \beta

In this framework, we model the implicit reward via an alternative parameterization, wherein the optimal policy coincides with $\pi_\theta$. Additionally, because our method corresponds to learning a reparameterized Bradley-Terry model, it inherits specific theoretical guarantees, notably consistency provided appropriate assumptions hold for the distribution of preference data \cite{bong2022generalized}. Further theoretical aspects of DPO and its connection to related literature are elaborated in Section~\ref{sec:theory}.

\textbf{Mechanistic interpretation of the DPO update.} To gain a detailed understanding of DPO’s operation, examining the gradient of the loss $\mathcal{L}_\text{DPO}$ proves insightful. The gradient with respect to the parameter vector $\theta$ can be expressed as:

\begin{multline*}\label{eq:gradient}
    \nabla_\theta \mathcal{L}_\text{DPO}(\pi_\theta;\pi_\text{ref}) = \\ -\beta\mathbb{E}_{(x, y_w, y_l) \sim \mathcal{D}} \left[\underbrace{\sigma(\hat{r}_\theta(x, y_l) - \hat{r}_\theta (x, y_w))}_\text{greater emphasis when reward estimates are misordered}\left[\underbrace{\nabla_\theta\log \pi(y_w \mid x)}_\text{promote $y_w$ occurrence} - \underbrace{\nabla_\theta\log\pi(y_l \mid x)}_\text{demote $y_l$ occurrence}\right]\right],
\end{multline*}

where $\hat{r}_\theta(x, y) = \beta \log \frac{\pi_\theta(y \mid x)}{\pi_\text{ref}(y \mid x)}$ defines the implicit reward in terms of the learned model $\pi_\theta$ relative to the reference $\pi_\text{ref}$ (see Section~\ref{sec:theory} for further discussion). Intuitively, this gradient acts to bolster the likelihood of the more favored responses $y_w$ while diminishing that of the less favored ones $y_l$. Critically, the gradient updates are weighted in proportion to how strongly the implicit reward model $\hat{r}_\theta$ incorrectly rates the less preferred completion higher, modulated by the factor $\beta$—in other words, reflecting both the error in ordering and the strictness of the KL regularization. Empirical results highlight the necessity of this weighting mechanism; omitting the coefficient leads to undesirable model behavior (see Appendix Table~\ref{tab:unlikelihood_generations}).

\textbf{DPO overview.} The standard DPO process consists of the following steps: 1) For each prompt $x$, generate completion samples $y_1, y_2 \sim \pi_\text{ref}(\cdot \mid x)$, annotate them using human preference labels, and build an offline preference dataset $\mathcal{D} = \{x^{(i)}, y_w^{(i)}, y_l^{(i)}\}_{i=1}^N$; 2) Train the language model $\pi_\theta$ by minimizing $\mathcal{L}_\text{DPO}$ with respect to the reference policy $\pi_\text{ref}$, the constructed dataset $\mathcal{D}$, and a target value of $\beta$. 
In applied settings, it is preferable to make use of existing publicly released preference datasets rather than producing new samples and collecting human feedback. Because these datasets are typically collected with $\pi^\text{SFT}$ as the base policy, $\pi_\text{ref}$ is set to $\pi^\text{SFT}$ wherever this model is available. In situations where $\pi^\text{SFT}$ cannot be used, $\pi_\text{ref}$ is instead chosen to maximize the likelihood of preferred completions $(x, y_w)$ according to ${\pi_\text{ref} = \argmax_{\pi}\mathbb{E}_{x, y_w \sim \mathcal{D}}\left[\log \pi(y_w \mid x)\right]}$. This method aims to address the potential mismatch between the ideal but unknown reference distribution and the practical $\pi_\text{ref}$ employed by DPO. A comprehensive account of implementation specifics and hyperparameter settings is provided in Appendix~\ref{app:implementation}.

\section{Theoretical Analysis of DPO} In what follows, we present deeper theoretical insights into DPO, clarify its theoretical underpinnings, and demonstrate how DPO addresses shortcomings in actor-critic methods such as those used in RLHF frameworks like PPO~\cite{schulman2017proximal}.

\label{sec:theory}

\subsection{Your Language Model Is Secretly a Reward Model} DPO achieves both the avoidance of training a separate reward model and circumventing reinforcement learning policy optimization by reducing the process to a single maximum likelihood objective. Importantly, the loss function in Eq. \ref{eq:main_eq} is functionally identical to a Bradley-Terry model parameterized with rewards as $r^*(x, y) = \beta \log\frac{\pi^*_\theta(y \mid x)}{\pi_\text{ref}(y \mid x)}$, and training the model $\pi_{\theta}$ becomes analogous to fitting a reward model per Eq. \ref{eq:reward_model} under an appropriate reparameterization. Here, we develop the theoretical justification for this change of variables, verify that it does not restrict the family of possible reward models, and prove that the optimal policy can be exactly retrieved. We commence by formalizing an equivalence relationship among reward functions.

\begin{definition}
Two reward functions $r(x, y)$ and $r'(x, y)$ are said to be equivalent if and only if ${r(x, y) - r'(x, y) = f(x)}$ for some function $f$.
\end{definition}

This equivalence is straightforward to verify as it indeed constitutes an equivalence relation, resulting in a partitioning of all reward functions into distinct equivalence classes. This leads to the following two lemmas:

\begin{lemma}\label{lemma:same_prefrence}
Within the Plackett-Luce, and specifically the Bradley-Terry, framework for modeling preferences, any two reward functions that are members of the same equivalence class will yield identical preference distributions.
\end{lemma}

\begin{lemma}\label{lemma:same_policy}
Reward functions from the same equivalence class will produce identical optimal policies for the corresponding constrained RL problem.
\end{lemma}

The proofs are elementary and thus postponed to Appendix \ref{app:lemma1}. The first lemma highlights a familiar under-identifiability feature of models belonging to the Plackett-Luce family \cite{plackett1975analysis}. This inherent ambiguity necessitates the adoption of additional identifiability criteria to secure meaningful guarantees for the maximum likelihood estimates in Eq. \ref{eq:reward_model} \cite{bong2022generalized}. The second lemma shows that the choice of reward function within an equivalence class does not affect the optimal policy, implying that our learning objective only needs to identify any member of the equivalence class yielding the optimal behavior. The following theorem is established in Appendix~\ref{app:thm1}:

\begin{theorem}\label{thm:main}
    Provided some mild regularity assumptions, every reward class that is compatible with Plackett-Luce (and particularly Bradley-Terry) models may be expressed using the reparameterization ${r(x, y) = \beta \log \frac{\pi(y\mid x)}{\pi_\text{ref}(y\mid x)}}$ for some model $\pi(y\mid x)$ and a chosen reference model $\pi_\text{ref}(y \mid x)$.
\end{theorem}

\begin{sproof}
    Let $r(x, y)$ be an arbitrary reward function, which determines its associated optimal model $\pi_r(y \mid x)$ as specified in Eq. \ref{eq:op_policy}. We aim to demonstrate that within the equivalence class of $r$, there exists a reward function of the stated reparameterized form. To this end, define the projection $f$ as
\begin{equation}
    f(r; \pi_\text{ref}, \beta)(x, y) = r(x, y) - \beta\log\sum_{y}\pi_\text{ref}(y\mid x)\exp\left(\frac{1}{\beta}r(x, y)\right)
\end{equation}
This mapping $f$ serves to adjust $r(x, y)$ by subtracting the log-partition function, depending solely on $x$, thus yielding another reward function in the same equivalence class. Substituting $r$ using the right-hand side of Eq.~\ref{eq:main_eq} (which applies universally to reward functions), we obtain $f(r; \pi_\text{ref}, \beta)(x, y) = \beta \log \frac{\pi_r(y\mid x)}{\pi_\text{ref}(y\mid x)}$. Hence, $f$ produces an equivalently classed reward function with the required structure, establishing that no expressivity is lost in adopting the proposed reparameterization.
\end{sproof}

Alternatively, Theorem~\ref{thm:main} can be interpreted as uniquely identifying, within each equivalence class, the particular reward function selected by the DPO reparameterization. Specifically, this reward function must satisfy:

\begin{equation}\label{eq:lag_p}
     \sum_{y}\underbrace{\pi_\text{ref}(y\mid x)\exp\left(\frac{1}{\beta}r(x, y)\right)}_{=\pi(y\mid x)\text{, using Thm.~\ref{thm:main} reparam.}} = 1,
\end{equation}

In other words, $\pi(y\mid x)$ constitutes a valid probability distribution: all probabilities are positive and sum to unity. Referring back to Eq.~\ref{eq:op_policy}, it becomes clear that Eq.~\ref{eq:lag_p} represents the partition function associated with the optimal policy determined by the reward function $r(x, y)$. The principal contribution of the DPO algorithm is the introduction of specific constraints on the otherwise under-constrained Plackett-Luce (and, in particular, Bradley-Terry) class of preference models. This constraint preserves the range of expressible reward models but crucially ensures that the optimal policy given by Eq.~\ref{eq:op_policy} becomes analytically manageable for every input prompt $x$.

\subsection{Instability of Actor-Critic Algorithms}  
Our framework additionally facilitates an analysis of the instabilities observed in standard actor-critic methods commonly employed in RLHF, such as PPO. We focus on the RL fine-tuning stage described in Section \ref{section:prelims}, adhering to the RLHF paradigm. By drawing on the control as inference perspective \cite{levine2018reinforcement}, we frame the constrained RL problem as outlined in \ref{eq:RL}. Assume a parameterized policy $\pi_{\theta}(y\mid x)$, and seek to minimize the divergence $\mathbb{D}_{\text{KL}}[\pi_{\theta}(y|x) \mid\mid \pi^*(y\mid x)]$, where $\pi^*$ denotes the optimal policy from Eq.~\ref{eq:optimum_model}, which arises from the reward function $r_{\phi}(y, x)$. Through algebraic manipulation, the resulting optimization objective can be expressed as:

\begin{equation}\label{eq:AC}
    \max_{\pi_{\theta}}\mathbb{E}_{\pi_{\theta}(y\mid x)}\left[\underbrace{r_{\phi}(x, y) -\beta\log\sum_{y}\pi_\text{ref}(y\mid x)\exp\left(\frac{1}{\beta}r_{\phi}(x, y)\right)}_{f(r_{\phi}, \pi_\text{ref}, \beta)} - \underbrace{\beta\log\frac{\pi_{\theta}(y\mid x)}{\pi_\text{ref}(y\mid x)}}_{\text{KL}}\right]
\end{equation}

The objective described here aligns with that targeted in earlier studies 
\citep{ziegler2020finetuning, stiennon2022learning, bai2022training, ouyang2022training}, wherein the DPO-equivalent reward function for the reward class $r_{\phi}$ is optimized. In this framework, the normalization component within $f(r_{\phi}, \pi_\text{ref}, \beta)$ can be viewed as representing the soft value function associated with the reference policy $\pi_\text{ref}$. Although this normalization factor does not change the nature of the optimal solution, its absence may lead to high variance in the policy gradient, thereby reducing the stability of the learning process. One approach is to estimate the normalization factor with a learned value function; however, this often proves challenging to train effectively. Previous work instead typically normalizes rewards with respect to a human completion baseline, effectively utilizing a single-sample Monte Carlo estimate for the normalization. By contrast, the DPO reparameterization produces a reward structure that inherently obviates the need for any such baselines.

\section{Experiments}
We proceed to experimentally assess the capability of DPO to train policies directly from preference data. First, within a controlled text-generation environment, we explore DPO's efficiency in balancing reward maximization against KL-divergence regularization relative to the reference policy, in comparison to standard preference-based methods such as PPO. Subsequently, we analyze DPO's effectiveness on tasks that are both larger in scale and more challenging, including applications in summarization and dialogue, using larger language models. Our results indicate that, with minimal hyperparameter adjustment, DPO frequently matches or outperforms robust baselines—including RLHF implementations using PPO and strategies that select the top sample among $N$ generated trajectories according to a learned reward model. We detail our experimental protocols prior to presenting results, with further specifics available in Appendix~\ref{app:exp_details}.

\textbf{Tasks.} The scope of our experiments spans three distinct open-ended text generation tasks. In each setting, the algorithms are tasked with acquiring a policy based on a dataset of preferences, $\mathcal{D}=\bigl\{x^{(i)}, y_w^{(i)}, y_l^{(i)}\bigr\}_{i=1}^N$. In the \textbf{controlled sentiment generation} scenario, $x$ represents the initial segment of a movie review sampled from the IMDb dataset \cite{maas-EtAl:2011:ACL-HLT2011}, and the objective is to generate a continuation $y$ that expresses positive sentiment. To enable a rigorous assessment, we synthesize preference pairs for generated outputs using a pre-trained sentiment classifier such that $p(\text{positive}\mid x,y_w)>p(\text{positive}\mid x,y_l)$. For SFT, the GPT-2-large model is further trained on positive reviews within the IMDB train set, until convergence is reached (additional procedural specifics appear in App~\ref{app:sentiment_details}). In the \textbf{summarization} task, $x$ corresponds to a Reddit forum post, and the goal is for the policy to generate a summary $y$ highlighting the main arguments of the post. In alignment with established literature, our work utilizes the Reddit TL;DR dataset \citep{volske-etal-2017-tl} along with the human preference annotations from \citeauthor{stiennon2022learning}. Our SFT baseline is obtained by additional fine-tuning on human-composed post summaries\footnote{\url{https://huggingface.co/CarperAI/openai_summarize_tldr_sft}} using the TRLX RLHF framework \citep{leandro_von_werra_2023_7790115}. Notably, the human preference labels originate from evaluations on outputs of another comparably trained SFT system. Lastly, for \textbf{single-turn dialogue}, $x$ is a human-generated prompt, potentially ranging from scientific inquiries to requests for advice. Here, the policy is tasked with producing a relevant and constructive reply $y$. For this purpose, we employ the Anthropic Helpful and Harmless dialogue dataset \citep{bai2022training}, which includes 170,000 annotated exchanges between a user and an automated assistant. Each dialogue terminates with two model-generated responses and an annotation marking the preferred option. As there is no off-the-shelf SFT model for this task, we create one by fine-tuning a generic language model exclusively on preferred responses.

\textbf{Evaluation.} We adopt two distinct evaluation methodologies in our experimental framework. To scrutinize the effectiveness of each algorithm with respect to the constrained reward maximization objective, we assess their performance in the controlled sentiment generation scenario by plotting their frontiers of realized reward versus KL-divergence relative to the reference policy. The feasibility of this evaluation arises from access to the actual reward function, which in this context is a sentiment classifier. By contrast, in practical applications where the true reward function is inaccessible, we rely on measuring the algorithms’ \textit{win rate} compared to a baseline policy. Here, GPT-4 serves as a stand-in for human evaluators, rating the quality of summaries and the helpfulness of responses in summarization and single-turn dialogue tasks, respectively. For the summarization experiments, the reference summaries from the test corpus act as the baseline, while in dialogue tasks, the dataset’s annotated preferred responses serve this role. Although recent research indicates that language models can outperform existing automated metrics as evaluators \citep{Chen2023ExploringTU}, we bolster the case for leveraging GPT-4 through an independent human study, the details of which are found in Sec.~\ref{sec:human-judgments}. Our analysis reveals that GPT-4's assessments closely align with those of humans, with human-GPT-4 agreement generally matching or exceeding the level of concordance observed among human annotators.

\textbf{Methods.} Beyond DPO, we benchmark several established techniques for training language models according to human preference data. In the summarization domain, we experiment with zero-shot prompting using \textbf{GPT-J} \citep{gpt-j}, and for the dialogue setting, we employ 2-shot prompting with \textbf{Pythia-2.8B} \citep{biderman2023pythia}. Our evaluation also encompasses the \textbf{SFT} approach and \textbf{Preferred-FT}, where the model is fine-tuned via supervised learning using the selected output $y_w$, sourced from the SFT model in sentiment and summarization, or from a general language model for single-turn dialogue. Another method in our suite is \textbf{Unlikelihood}~\citep{welleck2019neural}, which encourages the model to increase the likelihood of the preferred completion $y_w$ while actively decreasing the likelihood for the less desired completion $y_l$. The ‘unlikelihood’ loss component may be scaled by an optional coefficient $\alpha\in[0,1]$. We further include \textbf{PPO} \citep{schulman2017proximal}, which utilizes a reward function inferred from preference data, and \textbf{PPO-GT}, an oracle variant trained with access to the actual reward function (only possible in the controlled sentiment setting). In sentiment-focused experiments, we test two PPO-GT implementations: an out-of-the-box version \cite{leandro_von_werra_2023_7790115} and a custom variant featuring reward normalization and additional hyperparameter tuning—modifications that are also applied to standard PPO when using learned rewards. Lastly, we introduce the \textbf{Best of $N$} baseline: $N$ responses are sampled from either the SFT model or Preferred-FT (for dialogue), and the top response according to the learned reward model is chosen. While this method separates reward model quality from PPO’s optimization mechanics and achieves strong results, it becomes computationally costly for larger $N$, as it necessitates producing $N$ completions per query at inference time.

\subsection{How well can DPO optimize the RLHF objective?}

Typical RLHF methods optimize a reward function under a KL-divergence constraint, ensuring the new policy does not diverge excessively from a reference policy while still maximizing reward. Therefore, an appropriate comparison of such algorithms requires assessment of both achieved reward and KL divergence; improvements in reward are only meaningful if not accompanied by disproportionate increases in KL. Figure~\ref{fig:frontier-tldr-main} illustrates the tradeoff between reward and KL for several algorithms under the sentiment task. For a robust comparison, each method is run multiple times, varying the conservativeness parameter for the policy in each run: for PPO, we use target KL values $\in\{3,6,9,12\}$; for unlikelihood, $\beta \in \{0.05,0.1,1,5\}$ and $\alpha\in\{0.05,0.1,0.5,1\}$; and for preferred-FT, different random seeds. Across all methods and parameter choices, this results in a total of 22 training runs. At intervals of 100 training steps until convergence, every policy is evaluated using a designated set of test prompts, calculating both the mean reward based on the true reward function and the mean sequence-level KL\footnote{Specifically, this refers to the sum over per-timestep KL-divergences.} relative to the reference policy $\text{KL}\left(\pi\mid \mid \pi_\text{ref}\right)$. The data demonstrate that DPO defines a markedly superior frontier, attaining maximal reward while maintaining low KL divergence. This observation stands out for several reasons: despite sharing the same objective as PPO, DPO consistently exhibits greater efficiency, as its reward/KL tradeoff entirely supersedes that of PPO. Furthermore, DPO yields a more favorable frontier than PPO, \emph{even when PPO is granted access to the ground truth rewards} (PPO-GT).

\subsection{Can DPO scale to real preference datasets?}
\label{sec:dpo-real-datasets}
We proceed to assess the fine-tuning capabilities of DPO on tasks such as summarization and single-turn dialogue. In the context of summarization, traditional automatic metrics like ROUGE often exhibit weak correlation with human judgments~\citep{stiennon2022learning}, and previous research demonstrates that language models refined with PPO based on human feedback yield higher-quality summaries. Our experimental setup involves generating completions from various methods using the test split of the TL;DR summarization dataset, with the evaluation criterion being the mean win rate against test set reference completions. Samples are drawn at temperatures between 0.0 and 1.0, with the resulting win rates depicted in Figure~\ref{fig:frontier-tldr-main} (right). DPO, PPO, and Preferred-FT are all applied to the same GPT-J SFT model\footnote{\url{https://huggingface.co/CarperAI/openai_summarize_tldr_sft}} for fine-tuning. Notably, DPO achieves an approximate 61\% win rate at temperature 0.0, surpassing PPO’s best performance of about 57\% at the same temperature. DPO also reaches a higher peak win rate than the strongest $N$-best baseline. It is important to highlight that no extensive tuning was performed on DPO’s $\beta$ hyperparameter; as a result, the actual potential of DPO could be greater. In addition, DPO exhibits significantly greater resilience to variations in sampling temperature than PPO, whose performance may regress to that of the original GPT-J at higher temperatures. Preferred-FT yields only marginal improvements relative to the SFT baseline. We provide further direct comparison between DPO and PPO through human evaluation in Section~\ref{sec:human-judgments}, where human annotators favored DPO completions sampled at temperature 0.25 in 58\% of pairwise comparisons against PPO completions at temperature 0.

For single-turn dialogue, our experiments utilize a subset of the Anthropic HH test split \citep{bai2022training} that involves a single exchange between a human and an assistant. We assess methods using GPT-4-based evaluation, where model outputs are compared to the human-preferred completions from the test set to determine win rates. Due to the absence of a canonical SFT model for this particular benchmark, we begin with a pre-trained Pythia-2.8B model and employ Preferred-FT to fine-tune a reference model on the curated completions, ensuring the outputs remain close to the training distribution. Subsequently, we apply DPO training. As part of our analysis, we include a comparison to the strongest performance among 128 Preferred-FT completions (the Best of $N$ metric, which levels off at $N=128$ in our experiments; refer to Appendix Figure~\ref{fig:best-of-n}) and also to a two-shot prompted version of the original Pythia-2.8B model, observing that DPO matches or outperforms these baselines for optimal temperature selections. Furthermore, we analyze an RLHF model trained with PPO on the Anthropic HH dataset\footnote{\url{https://huggingface.co/reciprocate/ppo_hh_pythia-6B}}, leveraging implementations from a reputable repository\footnote{\url{https://github.com/CarperAI/trlx/tree/main/examples/hh}}. However, we are unable to identify a prompt configuration or temperature that yields superior results compared to the base Pythia-2.8B model. Drawing on our TL;DR findings and noting the shared reward function optimization, we treat the Best of 128 baseline as a practical proxy for the performance of PPO-trained models. Overall, DPO stands out as the only method in our study that both enhances efficiency and surpasses the preferred-completion baseline in the Anthropic HH dataset, while achieving outcomes on par with or exceeding the much more computationally intensive Best of 128 approach. Lastly, as shown in Figure~\ref{fig:dialogue-main}, DPO rapidly attains peak performance during training.

\subsection{Generalization to a new input distribution}

To more thoroughly assess how PPO and DPO perform under shifts in input distribution, we evaluate the policies trained via PPO and DPO on our Reddit TL;DR summarization task using a distinct dataset—specifically, the test split of news articles from CNN/DailyMail \citep{nallapati-etal-2016-abstractive}. For these evaluations, we use the best sampling temperatures identified for TL;DR (0 and 0.25). Table~\ref{tab:ood} presents these findings. The win rate of GPT-4 was computed in comparison to the gold-standard dataset summaries, employing the same GPT-4 (C) prompt as with Reddit TL;DR, except substituting ``news article'' for ``forum post''. On this new dataset, the DPO policy still notably exceeds the performance of PPO. These results offer preliminary evidence that DPO policies are able to generalize to novel input distributions comparably to, if not better than, PPO policies—even without the additional unlabeled Reddit TL;DR prompts utilized by PPO.

\subsection{Validating GPT-4 judgments with human judgments}
\label{sec:human-judgments}
We implement a human evaluation to assess the consistency of GPT-4’s judgments, drawing from results of the TL;DR summarization experiment and deploying two distinct prompts for GPT-4. The \textbf{GPT-4 (S)} (simple) prompt inquires simply about which summary better conveys the essential content of the post. Conversely, the \textbf{GPT-4 (C)} (concise) prompt includes an additional request for conciseness, prompted by our observation that with the \textbf{GPT-4 (S)} prompt, GPT-4 tends to prefer longer and more repetitive outputs than do human evaluators. The full text of both prompts is given in Appendix~\ref{app:prompts}. Our human study involves three experimental conditions—comparing the top-performing (DPO, temp. 0.25), lowest-performing (PPO, temp. 1.0), and a medium-performing (SFT, temp. 0.25) method, each against PPO with greedy decoding (the temperature value yielding its best results)—to capture a broad spectrum of output quality. Across both prompt formulations, we observe that GPT-4’s selections align with those of human evaluators at a rate comparable to inter-human agreement, suggesting GPT-4 is a robust surrogate for human judgment (note that, given the limited pool of human evaluators, multiple human assessments were collected only for the DPO and PPO-1 comparison). In general, win rates from the \textbf{GPT-4 (C)} prompt more closely mirror human preferences, motivating our use of this prompt for the primary analyses in Section~\ref{sec:dpo-real-datasets}. Further details about the human evaluation, including the rater interface and the list of volunteer participants, can be found in Appendix~\ref{app:human-study}.

\section{Discussion}
Preference-based learning serves as a robust and scalable approach for developing language models that are both capable and aligned with human intent. In this work, we presented DPO, a straightforward methodology for leveraging preferences in language model training without reliance on reinforcement learning. Rather than forcing preference optimization into conventional RL frameworks and employing standard RL algorithms, DPO establishes a direct correspondence between language model policies and reward functions. This allows language models to align with human preferences using a simple cross-entropy loss formulation, completely avoiding reinforcement learning while maintaining generality. Remarkably, DPO achieves competitive or superior results compared to prevailing RLHF methods, such as those based on PPO, and does so with minimal hyperparameter adjustment. As a consequence, DPO considerably lowers the entry barrier for preference-based training of language models.

\textbf{Limitations \& Future Work.} The findings in this study prompt a series of compelling directions for future research. One central question is how DPO-trained policies handle out-of-distribution scenarios relative to approaches trained via explicit reward functions. Preliminary evidence indicates that DPO-based models demonstrate comparable generalization to PPO-trained counterparts, but more systematic investigation is warranted. For instance, further work is needed to determine whether training through self-labeling from the DPO policy can also capitalize on unlabeled prompts. Another avenue of exploration involves understanding the manifestation of reward over-optimization within direct preference optimization—whether, for example, the moderate performance dip observed in Figure~\ref{fig:dialogue-main}-right is symptomatic of this phenomenon. Furthermore, our experiments were restricted to models containing up to 6B parameters; it remains to be seen how DPO scales to next-generation models with substantially larger parameter counts, a direction that holds considerable promise. Evaluation protocols also merit deeper analysis, as our results suggest that GPT-4-based win rates are sensitive to prompt formulation; optimizing automated evaluation techniques is an open problem. Beyond language modeling, DPO’s general principles may extend to preference-based training for generative models in diverse modalities, opening new applications for this paradigm.